% Source-level normalization and conditional local formula.
% GPT-6 Astra Pro; reasoning effort not exposed. 2026-11-21.
\subsection{SF-Q6-04: the actual stabilizer denominator and a full local expression}

\paragraph{Statement SF-Q6-04(a).}
Let $E$ be the host CM field and let $R$ be the rank datum used by
the genuine column construction.  For the denominator
\texttt{GlobalTheory.Mem.Col2.DGam E R} used in that construction,
every prime divisor satisfies
\[
 p\le [E:\mathbf Q]N_m+1.
\]
For $R=R_0$, this bound is $p\le641$.  In particular its
$863$-adic valuation is zero.

\begin{proof}
The source \texttt{PredMemGamRanks.lean} defines
\texttt{gamRanks R} as
\[
 \{N_j:1\le j\le m\}\cup
 \{\operatorname{rank}X',\operatorname{rank}Y_{m-1},n_V\}.
\]
Every rank in this set is at most $N_m$; the complete elementary
proof is also supplied as \texttt{le\_N\_of\_mem\_gamRanks} in
\texttt{DGamDvdOfMinkMult.lean}.
For each such rank $n$, \texttt{autOrd E n} is the set of orders
of the finite unitary automorphism groups of unimodular lattices
in $E^n$.  The source proves that the resulting set of orders is
finite and defines \texttt{DGam} to be their least common multiple.
It proves its positivity; the zero alternative in its definition
does not occur for the positive host involution.

Suppose a prime $p$ divides one of these finite group orders.
Cauchy's theorem gives an element $g$ of order $p$ in
$\mathrm{GL}_n(E)$.  Restriction of scalars realizes $g$ faithfully
on a rational vector space of dimension $[E:\mathbf Q]n$.
Its minimal polynomial divides $T^p-1$.  As $g\ne1$ and
$\mathbf Q$ has characteristic zero, the cyclotomic factor
$\Phi_p(T)$ divides that minimal polynomial.  Consequently
\[
 p-1=\deg\Phi_p\le [E:\mathbf Q]n\le[E:\mathbf Q]N_m.
\]
Every prime divisor of a finite least common multiple divides
one of its entries, which proves the assertion.  For $R_0$,
$[E:\mathbf Q]=2f=16$, $N_m=40$ and $16\cdot40+1=641$.

This is the actual denominator in the column chain:
\texttt{chainAt} sets \texttt{DG := DGam X.E R}.
The proof of \texttt{ColIn.Ok.chIn\_ok} uses this same denominator,
including the free initial stabilizer and all subsequent decorated
lattice stabilizers.  Thus the conclusion does not replace an
arbitrary denominator by a smaller one without justification.
\end{proof}

\paragraph{Statement SF-Q6-04(b): full expression at an inert prime.}
Let $D$ have the fixed ranks $(k,m,t,f)=(8,2,15,8)$ and let
$L:\texttt{SizeData}_2(D)$.  Fix a word $u$ and a coefficient
place $w$ above a rational prime $p>641$.  Assume the following
\emph{ties to the actual data}:
\begin{enumerate}
\item $p$ is neither a type prime $g_i$ nor the level prime $T_u$;
      \texttt{IsResidueClassZero D p}, $h_P(D,p)=0$, and
      $p\le X''$.
\item $L.F$ has the defining values \texttt{Fval}, $L.E$ has the
      defining values \texttt{Eval}, and the actual
      \texttt{ParametersFill} equality for the gain holds at $w$.
\item The chain denominator is the \texttt{DGam} of part (a),
      and the error constant is $c_{\mathrm{dig}}=3$.
\item At every CM embedding,
      $e_1=M+18$ and $e_2=M+54$.
\end{enumerate}
Then the full source expression, including the error term, is
\[
 \boxed{\quad
 \operatorname{sg}(u,w)
   =\frac{96M+3456}{p^2-1}
       -76800\log_p(a).
 \quad}
\]

\begin{proof}
The defining first branch of \texttt{gainP} is zero when
\texttt{IsResidueClassZero} and $h_P=0$ hold.  Hence
$\mathrm{gainD}_j=0$, whether or not $p$ lies in the drop set.
If $w$ is outside the support, the support axiom already gives
the same conclusion.  In the stated branch of \texttt{Fval},
\[
 F_j=\frac{(k-1)\sum_\tau e_{j,\tau}}{p^2-1}.
\]
The defining branch of \texttt{Eval} is
\[
 E_j=v_p(D_\Gamma)
       +c_{\mathrm{dig}}fN_m^2\log_p(a)
       +\frac{\sum_\tau e_{j,\tau}}{p^2-1}.
\]
Its stabilizer term is zero by (a).  There are $f=8$ CM
embeddings, so
$\sum_{j,\tau}e_{j,\tau}=8(2M+72)=16M+576$.
Subtracting the two errors from the two forced terms gives
\[
 (k-2)\frac{16M+576}{p^2-1}
       -2\cdot3\cdot8\cdot40^2\log_p(a),
\]
which is the asserted expression.
\end{proof}

\paragraph{A strictly positive numerical specialization of this expression.}
Set
\[
 p=863,\qquad x=10^{1000},\qquad
 \varepsilon=\frac1{100},\qquad J=1000,\qquad
 M=\lceil\exp(10^{10})\rceil.
\]
Writing $B=10^{10}$, the source scale gives
$b=\exp(B)$ and $a=10^{1000}\exp(B)$.
Since $M\ge B^2/2$, $\log 863>1$ and $\log10<3$, the displayed
full expression is strictly greater than
\[
 \frac{48B^2}{744768}-76800(B+3000)>0.
\]
The last assertion is an exact rational inequality, independently
checked by the accompanying script.  This uses the full
$F-E$ formula and the actual vanishing stabilizer term; it is not
a conclusion from positivity of $F$ alone.

\paragraph{Remaining realization condition.}
Part (b) is a conditional identity, and the final paragraph proves
positivity of its concrete right hand side.  It is not yet a
certificate that the complete actual global row, twist, integral
lattice and period-valuation data realizing all four ties have
been assembled.  In particular, the source's abstract
\texttt{RowsRestModel} completion of actual columns must not be
substituted for the actual Poincare map required by a
\texttt{RealPack}.

The starting full host $H_0$ has $(q,d)=(571,599)$,
$TT=\varnothing$, $T_2=\{0\}$, $L_w=50$ and split bound $2$.
It must first be relevelled: the complete proof of
\texttt{exists\_relevel\_primes\_one\_mod} preserves its raw host
field and these combinatorial fields while imposing the stronger
\texttt{Pre.level\_cond}.  Choosing its bound above $863$ makes
the new type primes and distinguished level primes exceed $863$.
The concrete scale has
$\lfloor x^{\varepsilon/(5J)}\rfloor
 =\lfloor10^{1/500}\rfloor=1\le2$,
so its split-bound condition does not require changing the raw host.
The remaining construction must be checked independently of the
\texttt{RankVal} route excluded by SF-Q6-03.

\paragraph{Exact source locators and checks.}
Paths are relative to \texttt{lean/src/}.
\begin{itemize}
\item \texttt{Development/AbelianVarietyTheory/Definitions/AutomorphicMembership/
PredMemGamRanks.lean}: definitions \texttt{autO}, \texttt{autOrd},
\texttt{gamRanks}, \texttt{gamOrd}, \texttt{DGam} and proofs
\texttt{autO\_finite} through \texttt{DGam\_pos} (lines 322--389);
\texttt{chainAt} (lines 955--966).
\item \texttt{Development/AbelianVarietyTheory/Proofs/IwasawaDecomposition/
DGamDvdOfMinkMult.lean}: \texttt{le\_N\_of\_mem\_gamRanks}
(lines 13--27); \texttt{TrueColumnChainsSatisfyConditions.lean}:
\texttt{ColIn.Ok.chIn\_ok} (lines 55--75).
\item \texttt{Formalism/LocalToGlobal/Columns/
DecorationStabilizerTrivialUnipotent.lean}:
\texttt{primeOrder\_le}, \texttt{prime\_le\_of\_dvd\_card}
(lines 345--389), the same restriction-of-scalars/Cauchy argument.
\item \texttt{Formalism/LocalRingCardinality/DataForced.lean}:
\texttt{FInert}, \texttt{FvalP} (lines 112--135);
\texttt{DataErr.lean}: \texttt{EatC} (lines 121--123),
\texttt{Eval} (lines 159--170); \texttt{DataParams$_2$.lean}:
\texttt{gainD}, \texttt{sg} (lines 212--216).
\item \texttt{Formalism/SharpRestriction/GrowthBounds/
IsResidueClassZero.lean}: \texttt{gainP} (line 355),
\texttt{ParametersFill} (lines 401--411).
\item \texttt{Development/AbelianVarietyTheory/Proofs/IwasawaDecomposition/
ExistsHostCongruencePrimes.lean}:
\texttt{exists\_relevel\_primes\_one\_mod} (lines 16--102).
\end{itemize}
The full cited stabilizer proofs and the complete defining local
branches were read and reconstructed.  No Lean compilation is claimed.
